% ============================================================================
\section{Conclusion and Next Steps}
\label{sec:conclusion}
% ============================================================================

\subsection{Summary}

This article has surveyed four paradigms for building neural networks
without traditional multiply-accumulate arithmetic:

\begin{enumerate}[leftmargin=2em]
  \item \textbf{Binary Neural Networks}~\cite{courbariaux2016bnn, rastegari2016xnornet} replace multiplication with
    XNOR and accumulation with popcount. They are the most mature
    paradigm~\cite{bnnsurvey2020}, but CPU performance depends on data
    packing and popcount overhead as well as arithmetic count
    ~\cite{fischer2025popcount}.

  \item \textbf{Tsetlin Machines}~\cite{granmo2018tsetlin} replace the neuron metaphor entirely
    with propositional logic clauses learned through simple automata.
    They are the only paradigm that avoids floating-point arithmetic
    during both training and inference, and they offer full
    interpretability of the learned model. They achieve 99.3\% on
    MNIST~\cite{lecun1998mnist} with the convolutional
    variant~\cite{granmo2019convtm}.

  \item \textbf{Logic Gate Networks}~\cite{petersen2022difflogic} learn both the type and
    connectivity of logic gates end-to-end. New connection-optimised
    training substantially reduces reported gate counts, making them a
    promising sparse Boolean building block rather than merely an
    accelerator-circuit technique~\cite{mommen2026fullytrainable}.

  \item \textbf{KAN/LUT compilation}~\cite{liu2024kan, lutkan2025} trains with maximally expressive
    learnable functions (B-splines on every edge) and then compiles
    them into lookup tables for deployment. This elegantly trades
    memory for computation, achieving near-zero-cost inference at
    the price of larger stored tables.
\end{enumerate}

\subsection{The Hybrid Hypothesis}

We propose a CPU-first search over composable building blocks: Tsetlin
clauses for interpretable sparse features, Logic Gate Networks for learned
Boolean composition, LUTs for small local nonlinearities, and ternary or
low-bit integer kernels where they beat an irregular Boolean graph. A
useful hybrid must achieve competitive accuracy on MNIST while:

\begin{itemize}[leftmargin=2em]
  \item Using \emph{no} floating-point arithmetic at inference time
  \item Producing a partially interpretable model
  \item Using packed CPU operations, local LUTs, and only the low-bit
    arithmetic that materially improves end-to-end runtime
  \item Being implementable in Rust with no external ML framework
    dependencies
\end{itemize}

\subsection{Scaling Beyond MNIST}

\Cref{sec:transformers-llms} and \cref{sec:scaling-llms} examined
whether the bitwise paradigms can scale beyond classification tasks.
The analysis reveals a nuanced picture:

\begin{itemize}[leftmargin=2em]
  \item \textbf{For LLMs:} No single paradigm can replace the
    transformer architecture, but each paradigm can serve as a
    component \emph{within} a transformer. The combination of BitNet
    ternary weights, LUT-compiled activations, Logic Gate Network
    MLP blocks, and Tsetlin Machine MoE routing could yield a
    transformer that is 95--99\% integer/bitwise, with floating point
    confined to the attention softmax and normalisation layers.
    Softmax-free architectures (Mamba, linear attention) offer a
    path to eliminating even that remaining floating-point dependency.

  \item \textbf{For world models and sound synthesis:} The hybrid
    architecture could serve as a fast, interpretable front-end for
    acoustic classification and parameter estimation, feeding into
    classical DSP backends. The KAN/LUT layer is particularly
    promising as a bridge between Boolean features and continuous
    synthesis parameters.
\end{itemize}

\subsection{Next Steps}

\begin{enumerate}[leftmargin=2em]
  \item \textbf{Theory document:} A detailed design document
    specifying the hybrid architecture's mathematical formulation,
    layer dimensions, training algorithm, and convergence properties.

  \item \textbf{Rust implementation:} Implement and benchmark a small
    menu of blocks independently: packed clause masks, SIMD-aware
    binary/ternary dot products, contiguous LGN layers, and cache-sized
    LUT activations. Compose only blocks that beat the same quantised
    integer baseline on the target CPU.

  \item \textbf{Benchmarking:} Systematic comparison against both an
    FP32 MLP and a small quantised integer MLP on accuracy, latency,
    model size, CPU cycles, cache misses, branch behaviour, and energy.

  \item \textbf{Analysis:} Inspection of learned Tsetlin clauses
    for interpretability, gate type distributions in the LGN layer,
    and LUT resolution sensitivity analysis.

  \item \textbf{LLM scaling exploration:} Staged research programme
    from BitNet + LUT activations through softmax-free attention
    to full integration (\cref{sec:path-forward-llms}).
\end{enumerate}

\subsection{A Broader Perspective}

The significance of bitwise neural computation extends beyond any
single task. If these approaches can match floating-point accuracy on
practical problems, they open the door to AI that runs on
microcontrollers, embedded sensors, and edge devices without
specialised hardware---truly democratising access to machine
intelligence.

The honest conclusion is that no single architecture dominates all
tasks. The hybrid approach's strength---everything is Boolean,
interpretable, and cache-resident---is also its constraint. But by
embedding bitwise components within proven architectures (like
transformers) rather than trying to replace those architectures
wholesale, the paradigms surveyed here gain a much broader reach.

The journey from MNIST to fully bitwise language models is long,
but each step along it produces a working, deployable, interpretable
system. That incremental usability may be the hybrid architecture's
most important property.
